% appendices/global_requirements.tex
% Apéndice A: Requisitos globales del sistema (frontend + backend + comunes)
% ============================================================================
%
% GUÍA: Este apéndice contendrá los requisitos del sistema COMPLETO,
%       incluyendo los de frontend (TFG de Alejandro Soler Sichling),
%       los de backend (ya recogidos en el \cref{ch:requirements}) y los
%       requisitos comunes.
%
%       Se incluye aquí para dar una visión global del proyecto coordinado.
%       El cuerpo del TFG (cap. 3) solo contiene los requisitos del backend;
%       este apéndice los contextualiza dentro del sistema completo.
%
%   Estructura propuesta:
%
%   \section{Requisitos comunes}
%   % TODO: Requisitos que afectan a ambos TFGs (frontend y backend).
%   %       P.ej.: protocolo de comunicación, formato de datos intercambiados,
%   %       modelo de autenticación compartido, etc.
%
%   \section{Requisitos del backend}
%   % Referencia al \cref{ch:requirements} para no duplicar:
%   % "Los requisitos funcionales y no funcionales del backend se recogen
%   %  en detalle en el \cref{ch:requirements} de esta memoria."
%
%   \section{Requisitos del frontend}
%   % TODO: Cuando Alejandro complete sus requisitos, incluirlos aquí.
%   %       Puede ser un \input{} de sus ficheros .tex o un resumen.
% ============================================================================

% TODO: Completar cuando el compañero de frontend aporte sus requisitos.

Este apéndice recoge los requisitos globales del \ac{scdee}, abarcando
tanto el subsistema \textit{frontend} como el \textit{backend} y los
requisitos comunes a ambos. Los requisitos específicos del \textit{backend}
se desarrollan en detalle en el \cref{ch:requirements}.

\section{Requisitos comunes}

% TODO: Añadir los requisitos compartidos entre frontend y backend.
% Ejemplos: contrato de API, formato de intercambio de datos, modelo de
% autenticación, etc.

[Pendiente de definición conjunta con el TFG de \textit{frontend}.]

\section{Requisitos del backend}

Los requisitos funcionales y no funcionales del \textit{backend} se
recogen en detalle en el \cref{ch:requirements} de esta memoria.

\section{Requisitos del frontend}

% TODO: Incluir los requisitos del frontend cuando estén disponibles.
% Opciones:
%   a) \input{appendices/rf_frontend} si se proporcionan ficheros .tex
%   b) Resumen textual si se proporcionan en otro formato

[Pendiente de aportación por parte de Alejandro Soler Sichling,
autor del TFG de \textit{frontend} del \ac{scdee}.]
